% Source: https://artofproblemsolving.com/community/c3h1648747_a_problem
% Stable identity: aops-topic-1648747-q01
The first three steps are valid: multiplying both sides of $x=y$ by $x$, subtracting $y^2$, and factoring are all valid operations, so $(x-y)(x+y)=y(x-y)$ is true. (Both sides equal $0$, since $x-y=0$.)

\textbf{Error 1.} The step from $(x-y)(x+y)=y(x-y)$ to $x+y=y$ divides both sides by $x-y$. Since $x=y$, this is division by $0$, which is not allowed. The equation $0\cdot(x+y)=0\cdot y$ holds for all $x,y$ and gives no information; for instance $x=y=1$ satisfies it but not $x+y=y$.

\textbf{Error 2.} Even if $x=0$ had been correctly derived, the next step divides by $x=0$ again, and $\dfrac{0}{0}$ is undefined, so $\dfrac {x}{x}=1$ fails.

So the argument breaks at $\boxed{\text{the division by } x-y=0}$ (and again at the division by $x=0$).
